\chapter{电场与引力场}\label{ch:g11:electric-gravitational-fields}

地球如何拉动一个它从未接触过的苹果？牛顿本人曾称
超距作用为“如此荒谬”，以至于他拒绝为之辩护。
现代的回答是：质量或电荷会改变其周围的空间，
——它产生\emph{场}——而置于其中的任何物体，
会对其所在位置的场作出反应。我们并排建立两个重要的例子，
直到类比成为形式上的同一。

\section{从超距作用到场}

\begin{definition}[场]\label{def:g11:electric-gravitational-fields:field}
\emph{场}\index{场}为空间中的每一点指定一个矢量。
源（电荷、质量）会在其周围各处产生场，
无论那里是否有东西能感受到它；试探物体则
只受到由场在\emph{其自身所在位置}处决定的力。
这一两步式图像取代了超距作用。
\end{definition}

\begin{remark}
场并非一种记账的技巧：它携带能量，源的变化
以有限速率经由场向外传播——光本身就是
一种传播中的场（见\cref{ch:g12:light-as-wave}）。这里只研究
不随时间变化的场。
\end{remark}

\section{电场}

\begin{definition}[电场]\label{def:g11:electric-gravitational-fields:efield}
若一小试探电荷 $q$ 置于点 $M$ 处，受到的电场力
为 $\vect F$，则 $M$ 点处的\emph{电场}\index{电场}为
\[
\vect E = \frac{\vect F}{q}
\qquad \text{（单位 \unit{N/C}，也写作 \unit{V/m}）}.
\]
$\vect E$ 仅取决于场源及点 $M$，与试探电荷无关：
将 $q$ 加倍时 $\vect F$ 也加倍，商保持不变。
\end{definition}

\begin{proposition}[电荷所受的力]\label{prop:g11:electric-gravitational-fields:force}
电荷 $q$ 置于场 $\vect E$ 中，所受的力
$\vect F = q \vect E$：沿 $\vect E$ 方向（$q > 0$ 时），或与 $\vect E$ 相反（
$q < 0$ 时），大小为 $F = \abs{q}\,E$。
\end{proposition}

\begin{proof}
对定义式移项即可；符号规律来自负数乘法的代数运算。
\end{proof}

\begin{example}[读取场]\label{ex:g11:electric-gravitational-fields:reading}
一个试探电荷 $q = \qty{+2.0e-8}{C}$ 位于 $M$ 点，受到的力为
\qty{6.0e-4}{N}，方向向东。$M$ 点处的场为
$E = \num{6.0e-4} / \num{2.0e-8} = \qty{3.0e4}{N/C}$，方向向东；
$q' = \qty{-1.0e-8}{C}$ 的电荷置于同一点，受到
的力 $F = \num{1.0e-8} \times \num{3.0e4} = \qty{3.0e-4}{N}$——方向向西。
\end{example}

\begin{proposition}[点电荷的场]\label{prop:g11:electric-gravitational-fields:point}
点电荷 $Q$ 在距离 $d$ 处产生一个场，场的方向沿
电荷与该点的连线：背离 $Q$（$Q > 0$ 时），指向它（
$Q < 0$ 时）——大小为
\[
E = k\,\frac{\abs{Q}}{d^2},
\qquad k = \qty{8.99e9}{N.m^2/C^2}.
\]
\end{proposition}

\begin{proof}
由库仑定律（见\cref{ch:g11:fundamental-interactions}）可得
试探电荷 $q$ 在距离 $d$ 处所受的力：大小为 $k \abs{qQ} / d^2$，
沿连线方向，同号时为斥力。两边同除以 $q$，并
应用\cref{prop:g11:electric-gravitational-fields:force}。
\end{proof}

\begin{example}[带电球体]\label{ex:g11:electric-gravitational-fields:sphere}
一个范德格拉夫起电机的球体带有 $Q = \qty{+5.0e-7}{C}$（小带电球体
的作用如同位于其中心的点电荷）。当 $d = \qty{0.50}{m}$ 时：
\[
E = \num{8.99e9} \times \num{5.0e-7} / 0.50^2 \approx \qty{1.8e4}{N/C},
\]
方向沿半径背离球心。在 \qty{1.0}{m} 处，场降为
原来的四分之一：力中的 $1/d^2$ 规律在场中依然成立。
\end{example}

\begin{definition}[场线]\label{def:g11:electric-gravitational-fields:lines}
称\emph{场线}\index{场线}为处处与场相切的曲线，
即箭头所示的场方向。规则有三条：
\begin{itemize}
  \item 场线从正电荷出发，止于负电荷；
  \item 场线密集处场较强，稀疏处较弱；
  \item 两条场线永不相交：场在每一点只有一个方向。
\end{itemize}
\end{definition}

\begin{omfigure}
\begin{tikzpicture}[scale=0.82]
  \fill[omDef!20] (0,0) circle (0.32);
  \draw[omDef, thick] (0,0) circle (0.32);
  \node[font=\small] at (0,0) {$+$};
  \foreach \a in {0,45,...,315}
    \draw[-Stealth, omThm, thick] (\a:0.32) -- (\a:1.5);
\end{tikzpicture}\qquad\qquad
\begin{tikzpicture}[scale=0.82]
  \fill[omProp!20] (0,0) circle (0.32);
  \draw[omProp, thick] (0,0) circle (0.32);
  \node[font=\small] at (0,0) {$-$};
  \foreach \a in {0,45,...,315}
    \draw[-Stealth, omThm, thick] (\a:1.5) -- (\a:0.32);
\end{tikzpicture}

{\small 点电荷的场线：从 $+$ 径向射出，射入 $-$；
靠近电荷处场线密集，描绘了场强按 $1/d^2$ 增大的情形。}
\end{omfigure}

\section{电容器的匀强电场}

\begin{definition}[平行板电容器]\label{def:g11:electric-gravitational-fields:capacitor}
两块正对、平行的金属板，相互靠近，带有
异号电荷，构成\emph{电容器}\index{电容器}。两板之间
的电场为\emph{匀强电场}\index{匀强电场}：大小和
方向处处相同——与极板垂直，由正极板指向
负极板。
\end{definition}

\begin{proposition}[两板间的电场]\label{prop:g11:electric-gravitational-fields:uniform}
若两板间的电压（\cref{ch:g11:circuits-and-power}）为
$U$，两板间距为 $d$，则两板间电场的大小为
\[
E = \frac{U}{d} \qquad \text{（单位 \unit{V/m}）}.
\]
\end{proposition}
\admitted

\begin{remark}
匀强性与 $E = U/d$ 在本科一年级分册中予以严格推导；
此式解释了单位 \unit{V/m}。间隙很窄时，电场很强：
一毫米间隙上加 \qty{100}{V} 的电压，场强已达 \qty{e5}{V/m}。
\end{remark}

\begin{example}[电子束的偏转]\label{ex:g11:electric-gravitational-fields:beam}
在示波器中，电子束穿过极板，两板间
电压为 $U = \qty{100}{V}$，间距为 $d = \qty{5.0}{mm}$：
$E = 100 / \num{5.0e-3} = \qty{2.0e4}{V/m}$，每个电子受到
$F = \num{1.60e-19} \times \num{2.0e4} = \qty{3.2e-15}{N}$ 的力，方向指向
正极板——是电子重力 \qty{8.9e-30}{N} 的 $\num{4e14}$ 倍，
因此电子束会明显偏转（曲线轨迹
属于\cref{ch:g12:newtons-laws}的内容）。
\end{example}

\begin{omfigure}
\begin{tikzpicture}[scale=0.9]
  \draw[omDef, very thick] (0,2.2) -- (5.2,2.2)
    node[right, font=\small] {$+$};
  \draw[omProp, very thick] (0,0) -- (5.2,0)
    node[right, font=\small] {$-$};
  \draw[-Stealth, omThm] (0.8,2.05) -- (0.8,0.15);
  \draw[-Stealth, omThm] (1.8,2.05) -- (1.8,0.15);
  \draw[-Stealth, omThm] (2.8,2.05) -- (2.8,0.15)
    node[pos=0.25, right, font=\small] {$\vect E$};
  \draw[-Stealth, omThm] (4.6,2.05) -- (4.6,0.15);
  \fill[omDef] (3.7,1.4) circle (2.5pt)
    node[right=2pt, font=\small] {$q>0$};
  \draw[-Stealth, omMeth, very thick] (3.7,1.4) -- (3.7,0.6)
    node[left, font=\small] {$q\vect E$};
  \draw[Stealth-Stealth, thin] (-0.35,0) -- (-0.35,2.2)
    node[midway, left, font=\small] {$d$};
\end{tikzpicture}

{\small 电容器的匀强电场：平行、等间距的
场线，由 $+$ 极板指向 $-$ 极板，$E = U/d$；正电荷
沿场方向受力，负电荷则逆场方向受力。}
\end{omfigure}

\section{引力场}

\begin{definition}[引力场]\label{def:g11:electric-gravitational-fields:gfield}
若将检验质量 $m$ 置于某点，它所受的万有引力（即
重力）$\vect P$，则该处的\emph{引力场}\index{引力场}
为
\[
\vect g = \frac{\vect P}{m} \qquad \text{（单位 \unit{N/kg}）},
\]
反之亦然，$\vect P = m \vect g$。前一章已引入
数值 $g$；$\vect g$ 则增加了下落石块开始运动的方向。
\end{definition}

\begin{proposition}[地球的引力场]\label{prop:g11:electric-gravitational-fields:earth}
在距地心为 $d$ 处（地球质量为 $M$），
引力场指向地心，大小为
\[
g = G\,\frac{M}{d^2},
\qquad G = \qty{6.67e-11}{N.m^2/kg^2}.
\]
\end{proposition}

\begin{proof}
由万有引力（\cref{ch:g10:universal-gravitation}）可得
检验质量 $m$ 所受的力：大小为 $G M m / d^2$，方向指向中心；
两边除以 $m$。在地表处 $d = R$，此即
前一章中的 $g = GM/R^2 \approx \qty{9.81}{N/kg}$。
\end{proof}

\begin{example}[月球所在处的引力场]\label{ex:g11:electric-gravitational-fields:moon}
在月地距离 $d = \qty{3.84e8}{m}$ 处：
\[
g = \num{6.67e-11} \times \num{5.97e24} / (\num{3.84e8})^2
\approx \qty{2.7e-3}{N/kg}.
\]
地球的引力场永不消失，只是按 $1/d^2$ 减弱；这正是
使月球保持在轨道上运行的引力场。
\end{example}

\begin{remark}[近地表的匀强性]\label{rem:g11:electric-gravitational-fields:uniformg}
在实验室乃至山地的尺度上，$d$ 与
$R = \qty{6.37e6}{m}$ 相比几乎不变，因此 $\vect g$ 在很高精度下是匀强的：
竖直的场线，大小恒为 \qty{9.81}{N/kg}——地球的引力场看起来
与电容器的电场相似，正如站在花园里看地球，它也显得是平的。
\end{remark}

\begin{omfigure}
\begin{tikzpicture}[scale=0.9]
  \fill[omExo!15] (0,0) circle (0.85);
  \draw[omExo, thick] (0,0) circle (0.85);
  \node[font=\small] at (0,0) {地球};
  \foreach \a in {0,60,...,300}
    \draw[-Stealth, omThm, thick] (\a:2.1) -- (\a:0.95);
  \node[omThm, font=\small] at (30:1.9) {$\vect g$};
\end{tikzpicture}

{\small 地球的引力场：沿径向，指向
中心，大小为 $GM/d^2$——正是
负点电荷电场的精确图像。}
\end{omfigure}

\section{一个思想，两种力}

本章的两个场在数学上是同一套内容，只是换了两种外衣：

\begin{center}
\begin{tabular}{lcc}
 & \textbf{电场} & \textbf{引力场} \\
\hline
源 & 电荷 $Q$ & 质量 $M$ \\
场 & $\vect E = \vect F / q$ & $\vect g = \vect P / m$ \\
单位 & \unit{N/C} = \unit{V/m} & \unit{N/kg} \\
点源 & $E = k \abs{Q} / d^2$ & $g = G M / d^2$ \\
常量 & $k = \qty{8.99e9}{N.m^2/C^2}$ & $G = \qty{6.67e-11}{N.m^2/kg^2}$ \\
试探体受力 & $\vect F = q \vect E$ & $\vect P = m \vect g$ \\
场线 & 从 $+$ 出发，止于 $-$ & 总是指向质量 \\
\end{tabular}
\end{center}

对于左边的每个公式，
在对应关系 $q \leftrightarrow m$、$\vect E \leftrightarrow \vect g$、$k \leftrightarrow G$ 下，都变为右边的公式。

\begin{remark}[类比失效之处]\label{rem:g11:electric-gravitational-fields:noanalogue}
电荷有正、负两种，质量只有一种。因此电场可以
被抵消：金属盒内重新分布的电荷会消除盒内的任何外电场。
（即\emph{法拉第笼}——这就是为什么电梯会切断
手机信号），而正负电荷（$+$ 与 $-$）平衡的大块物质
在电学上是“沉默”的。引力既不能被屏蔽，也不能被中和：
每一千克质量都会增加它的拉力——这正是为什么该力比电力弱
$\num{e39}$ 倍（见\cref{exo:g11:electric-gravitational-fields:10}），
它才在大尺度上主宰宇宙。
\end{remark}

\section{练习}

\begin{exercise}[$\star$]\label{exo:g11:electric-gravitational-fields:1}
一个试探电荷 $q = \qty{+2.0e-6}{C}$ 位于点 $M$，受到
大小为 \qty{5.0e-3}{N}、方向朝北的力。求 $M$ 点的场（大小与方向），
再求电荷 $q' = \qty{-4.0e-6}{C}$ 放在 $M$ 点时所受的力。
\end{exercise}

\begin{exercise}[$\star$]\label{exo:g11:electric-gravitational-fields:2}
点电荷 $Q = \qty{+5.0}{nC}$ 位于 $O$ 点。
\begin{enumerate}
  \item 求距离 $d = \qty{30}{cm}$ 处的场：大小与方向。
  \item 场的大小减为一半时，距离为多少？
\end{enumerate}
\end{exercise}

\begin{exercise}[$\star$]\label{exo:g11:electric-gravitational-fields:3}
一个电容器两极板间的电压为 $U = \qty{600}{V}$，两板相距 \qty{3.0}{mm}。
求该处的场，再求该处一个电子所受的力。
($e = \qty{1.60e-19}{C}$).
\end{exercise}

\begin{exercise}[$\star$]\label{exo:g11:electric-gravitational-fields:4}
在火星上，质量为 \qty{250}{kg} 的探测器重 \qty{930}{N}。求火星
表面的场 $g_{\text{火星}}$，再求质量为 \qty{80}{kg}
的宇航员在那里的重力。
\end{exercise}

\begin{exercise}[$\star$]\label{exo:g11:electric-gravitational-fields:5}
判断下列说法是否正确，并各说明一条理由：(a) 两条场线
可以在场强较大处相交；(b) 场线从负电荷出发，并
止于正电荷；(c) 场线密集处表示场强较大。
\end{exercise}

\begin{exercise}[$\star\star$]\label{exo:g11:electric-gravitational-fields:6}
国际空间站在离地高度 \qty{420}{km} 处绕地球运行
($R = \qty{6.37e6}{m}$, $M = \qty{5.97e24}{kg}$).
\begin{enumerate}
  \item 求空间站处的 $g$；它是地表值的几分之几？
  \item 宇航员处于漂浮状态。请用一句话
        结合你的上一问结果解释这一现象。
\end{enumerate}
\end{exercise}

\begin{exercise}[$\star\star$]\label{exo:g11:electric-gravitational-fields:7}
多大的场强才能与电子的重力平衡
（$m_e = \qty{9.11e-31}{kg}$）？将其与
\cref{ex:g11:electric-gravitational-fields:beam}中的 \qty{2.0e4}{V/m} 比较，并由此得出
引力在电子物理中的地位。
\end{exercise}

\begin{exercise}[$\star\star$]\label{exo:g11:electric-gravitational-fields:8}
一颗质量为 \qty{2.0e-6}{kg} 的尘埃颗粒，在水平极板之间静止漂浮，
两板相距 \qty{2.0}{cm}，两板间电压为 $U = \qty{5.0}{kV}$。
\begin{enumerate}
  \item 求场强，再求颗粒所带电荷量的大小。
  \item 颗粒带正电，哪一块极板是正极板？
  \item 颗粒带有多少个元电荷？
\end{enumerate}
\end{exercise}

\begin{exercise}[$\star\star$]\label{exo:g11:electric-gravitational-fields:9}
电荷 $+Q$ 与 $+4Q$ 相距 \qty{30}{cm}。在连接二者的线段上，
何处的合场强为零？请先解释为什么该点必须
位于两电荷之间，且更靠近电荷量较小的那一个。
\end{exercise}

\begin{exercise}[$\star\star$]\label{exo:g11:electric-gravitational-fields:10}
在氢原子中，质子（$m_p = \qty{1.67e-27}{kg}$，电荷
$+e$）与电子（$m_e = \qty{9.11e-31}{kg}$，电荷 $-e$）
相距 $d = \qty{5.3e-11}{m}$。求二者之间的电力和
万有引力，并求它们的比值。
\end{exercise}

\begin{exercise}[$\star\star$]\label{exo:g11:electric-gravitational-fields:11}
在点 $M$ 处，若只有一个场源单独存在，它产生的电场大小为 \qty{300}{N/C}，
方向向东；若只有第二个场源单独存在，它产生的电场大小为 \qty{400}{N/C}，方向向北。电场按矢量叠加，
求点 $M$ 处的合电场（写出大小和方向），然后
求电荷 $q = \qty{-2.0e-6}{C}$ 置于该点时所受的电场力。
\end{exercise}

\begin{exercise}[$\star\star\star$]\label{exo:g11:electric-gravitational-fields:12}
地球的引力场在何高度处降为地表值的一半？
在高度 $h = R$ 处，引力场为地表值的几分之几？
\end{exercise}

\begin{exercise}[$\star\star\star$]\label{exo:g11:electric-gravitational-fields:13}
两个电荷量均为 $\qty{+2.0}{nC}$ 的点电荷位于 $A$、$B$ 两点，相距 \qty{6.0}{cm}。
点 $M$ 位于 $[AB]$ 的中垂线上，
距其中点 \qty{4.0}{cm}。
\begin{enumerate}
  \item 求距离 $AM$，再求每个电荷在 $M$ 处
        产生的电场大小。
  \item 利用对称性，将两个电场按矢量相加，
        求 $M$ 处合电场的大小和方向。
\end{enumerate}
\end{exercise}

\begin{exercise}[$\star\star\star$]\label{exo:g11:electric-gravitational-fields:14}
试利用电荷的两种正负，说明金属盒中的自由电荷如何
使盒内各处的外电场被抵消——并说明为什么
任何质量的排布都无法对引力做到同样的事。
引力屏蔽需要哪种不可能存在的东西？
\end{exercise}

\begin{exercise}[$\star\star\star$]\label{exo:g11:electric-gravitational-fields:15}
一个小球（$m = \qty{0.50}{g}$，$q = \qty{+2.0e-7}{C}$）用细线悬挂在
水平匀强电场 $E = \qty{1.0e4}{V/m}$ 中。小球静止时，
细线与竖直方向的夹角为 $\theta$。
\begin{enumerate}
  \item 列出小球受到的三个力；说明为什么 $\tan\theta = qE / (mg)$。
  \item 求 $\theta$，再求细线的张力。
\end{enumerate}
\end{exercise}

\section{问题：密立根的油滴}

\begin{problem}[{周末问题——称量电子的电荷：两极板间悬停的
油滴表明，电荷以不可分割的台阶出现，
每个台阶为 $\num{1.60e-19}$ 库仑}]
\label{pb:g11:electric-gravitational-fields:1}
1909 年，密立根在水平放置的
电容器两极板之间喷入油滴。调节电压，使选定的油滴
悬停不动——电场力恰与重力平衡——他便能
称出其电荷。本题中油滴质量为 $m = \qty{4.90e-15}{kg}$
（关闭电场时由其缓慢下落测得，此处视为已知）；
两极板间距 $d = \qty{8.0}{mm}$，\emph{上}极板
带正电，油滴所带电荷为负。

\medskip
\noindent\textbf{第一部分——装置。}
\begin{enumerate}
  \item 描述两极板间的电场；“匀强”有什么好处？
  \item 当 $U = \qty{800}{V}$ 时，求 $E$。
  \item 油滴所受电场力指向何方？为何
        \emph{上}极板必须带正电？
  \item 求油滴的重力。
  \item 将静止时的平衡条件写成
        两个力大小相等的等式。
\end{enumerate}

\medskip
\noindent\textbf{第二部分——平衡。}
\begin{enumerate}[resume]
  \item 由此式，将 $\abs{q}$ 用 $m$、$g$、$d$ 和 $U$ 表示。
  \item 当 $U = \qty{800}{V}$ 时油滴静止悬挂，求
        $\abs{q}$。
  \item 将结果除以 $e = \qty{1.60e-19}{C}$，油滴带有
        多少个多余电子？
  \item 将电压微调至 \qty{900}{V}，此时合力
        指向何方？
  \item 证明平衡电压为 $U = mgd/\abs{q}$：
        电荷越大，平衡电压越小。
\end{enumerate}

\medskip
\noindent\textbf{第三部分——台阶。}
一束辐射可使电子落到油滴上或从油滴上脱离；每次
电荷变化后，密立根都重新调节电压使其平衡。
\begin{enumerate}[resume]
  \item 经一次辐射后，平衡需要 $U = \qty{1200}{V}$，求
        新的 $\abs{q}$。油滴是得到还是失去了一个电子？
  \item 依次得到的平衡电压为 $U = 1200$、$800$、$600$、
        $480$ 和 \qty{400}{V}，求这五个电荷量。
  \item 证明这五个电荷量都是同一量的整数倍，并给出
        该量的值。
  \item 相邻两次平衡之间，电荷量
        变化多少？并作出解释。
  \item 解释为何此油滴在
        $U = \qty{700}{V}$ 时永远无法静止悬挂，无论怎样耐心调节。
  \item 密立根对数百个油滴重复了这一实验：每次测得的
        电荷量都是同一数值的整数倍。写出结论，
        并指出测得的是哪个量。
\end{enumerate}

\medskip
\noindent\textbf{第四部分——为什么没有引力版的密立根实验。}
\begin{enumerate}[resume]
  \item 在本章的对照表中，何者分别扮演 $q$、$E$
        以及 $qE$ 在第二部分平衡中的角色？
  \item 若改由上方的质量托住油滴，是否可行？求
        质量为 \qty{1000}{kg} 的铅球位于油滴上方 \qty{1.0}{m} 处时所产生的引力场，
        油滴所受的力，并与其重力相比较。
  \item 平衡需要一个足以对抗整个地球的向上拉力，
        为什么一块带电的极板能够提供，而任何质量的
        排布都做不到？
  \item 油滴的\emph{质量}也以台阶的形式变化（一个油分子
        约为 \qty{5e-25}{kg}）。比较一个电子落上时
        $m$ 与 $\abs{q}$ 的相对跳变，并作出结论：为何
        此平衡能够探测到电的“原子”，却探测不到
        油的分子？要点：电荷是量子化的，其“原子”
        即 $e = \qty{1.60e-19}{C}$。
\end{enumerate}
\end{problem}
